% =============================================================
% บทที่ 5 - บทสรุป
% แก้ไขเนื้อหาในบทนี้ได้ที่ไฟล์ chapter5.tex
% =============================================================

% -----------------------------------------------------------
\section{สรุปผลการดำเนินโครงงาน}
% -----------------------------------------------------------

โครงงานนี้ได้พัฒนาระบบวิเคราะห์และทำนายผลการแข่งขันฟุตบอลยุโรปด้วยกระบวนการวิทยาการข้อมูลเชิงตัวเลข (Data Science Pipeline) และแบบจำลองการเรียนรู้ของเครื่อง ซึ่งมีวัตถุประสงค์เพื่อประมวลผลตัวแปรบริบทก่อนแข่ง ค้นหาการเดิมพันที่มีมูลค่าคาดหมายเป็นบวก (+EV) และประเมินผลตอบแทนทางการเงินผ่านการทดสอบย้อนหลัง ผลการดำเนินงานสรุปเทียบกับวัตถุประสงค์ได้ดังตารางที่ \ref{tab:obj_result}

\begin{table}[htbp]
\centering
\caption{สรุปผลเทียบกับวัตถุประสงค์ของโครงงาน}
\label{tab:obj_result}
\begin{tabular}{|c|p{4.5cm}|p{5.5cm}|c|}
\hline
\textbf{ข้อ} & \textbf{วัตถุประสงค์} & \textbf{ผลที่ได้} & \textbf{สถานะ} \\
\hline
1 & เพื่อพัฒนากระบวนการวิทยาการข้อมูลแบบ end-to-end จากหลายแหล่งข้อมูล & สร้างท่อส่งข้อมูลโมดูลาร์ รวบรวมข้อมูลจาก Understat, Transfermarkt และ Football-Data.co.uk พร้อมระบบตรวจสอบคุณภาพข้อมูลด้วย Pandera & บรรลุ \\ \hline
2 & เพื่อวิเคราะห์ตัวแปรก่อนแข่ง เช่น Elo, ดัชนีความเหนื่อยล้า และฟอร์มผู้เล่น 11 ตัวจริง & ทำการทดสอบทางสถิติและ Feature Importance พบว่า Elo Rating และฟอร์มผู้เล่นมีอิทธิพลต่อแบบจำลองสูงสุดอย่างมีนัยสำคัญ & บรรลุ \\ \hline
3 & เพื่อพัฒนาและประเมินแบบจำลองทำนายความน่าจะเป็นของผลการแข่งขัน & พัฒนาแบบจำลอง Dixon-Coles และ HistGradientBoosting โดยทำค่า Log Loss ได้ $0.9741$ และ Accuracy $54.32\%$ สูงกว่า Baseline & บรรลุ \\ \hline
4 & เพื่อทดสอบย้อนหลังทางการเงินและประเมินกลยุทธ์การบริหารเงินทุน (+EV) & พัฒนาระบบ Backtesting จำลองการเดิมพัน 10 ฤดูกาล พบว่ากลยุทธ์ Fractional Kelly ให้ ROI $+5.42\%$ และควบคุม Drawdown ได้ดี & บรรลุ \\
\hline
\end{tabular}
\end{table}

จากตารางที่ \ref{tab:obj_result} โครงงานนี้บรรลุวัตถุประสงค์ทั้ง 4 ข้ออย่างสมบูรณ์

% -----------------------------------------------------------
\section{ข้อจำกัดของระบบ}
% -----------------------------------------------------------

\subsection{ข้อจำกัดด้านการทำงาน}

\begin{itemize}
    \item \textbf{ความพร้อมของข้อมูลไลน์อัป (Lineup Availability):} รายชื่อผู้เล่น 11 ตัวจริงต้องรอการประกาศอย่างเป็นทางการล่วงหน้าก่อนแข่งประมาณ 60 นาที ทำให้แบบจำลองไม่สามารถประเมินผลในระยะยาวล่วงหน้าหลายวันได้
    \item \textbf{ความผันผวนของราคาอัตราต่อรอง (Odds Movement):} ระบบยังไม่ได้บันทึกความเคลื่อนไหวของราคาต่อรองแบบเรียลไทม์ (Live Odds) แต่ใช้อัตราต่อรองย้อนหลังคงที่จาก Football-Data.co.uk
\end{itemize}

\subsection{ข้อจำกัดด้านฮาร์ดแวร์และซอฟต์แวร์}

\begin{itemize}
    \item \textbf{หน่วยความจำและการประมวลผล:} การคำนวณตัวแปรระดับผู้เล่น (Player-level features) ย้อนหลัง 10 ฤดูกาลครอบคลุม 5 ลีกยุโรป ต้องใช้หน่วยความจำ RAM ขั้นต่ำ 8 GB และระบบประมวลผลแบบ Multi-core เพื่อให้รัน Data Pipeline ได้มีประสิทธิภาพ
    \item \textbf{ข้อจำกัดการเข้าถึง API แหล่งข้อมูล:} การดึงข้อมูลสถิติจากเว็บไซต์ภายนอกขึ้นอยู่กับโครงสร้าง HTML และขีดจำกัดการรับส่งข้อมูล (Rate Limit) ของเว็บไซต์ต้นทาง
\end{itemize}

% -----------------------------------------------------------
\section{ปัญหาอุปสรรค และแนวทางแก้ไข}
% -----------------------------------------------------------

ปัญหาที่พบระหว่างดำเนินงานและแนวทางแก้ไขแสดงดังตารางที่ \ref{tab:problems}

\begin{table}[htbp]
\centering
\caption{ปัญหาอุปสรรคและแนวทางแก้ไขในการดำเนินโครงงาน}
\label{tab:problems}
\begin{tabular}{|c|p{5cm}|p{6.5cm}|}
\hline
\textbf{ลำดับ} & \textbf{ปัญหาที่พบ} & \textbf{แนวทางแก้ไข} \\
\hline
1 & ชื่อสโมสรและนักเตะใน Understat, Transfermarkt และ Football-Data ไม่ตรงกัน & จัดทำระบบ Alias Table เพื่อทำ Entity Resolution แปลงชื่อให้เป็นมาตรฐานเดียวกันก่อนเชื่อมโยงข้อมูล \\ \hline
2 & รูปแบบและชนิดข้อมูลเกิดความคลาดเคลื่อนระหว่างขั้นตอนใน Pipeline & กำหนดสเปกข้อมูลด้วย Pandera Schema Validation ตรวจสอบโครงสร้างข้อมูลทุกสเตปก่อนส่งต่อไปฝึกโมเดล \\ \hline
3 & ปัญหาข้อมูลอนาคตรั่วไหล (Data Leakage) หากใช้สถิติหลังจบการแข่งขัน & ออกแบบระบบให้รองรับเฉพาะข้อมูลก่อนเริ่มแข่ง (Pre-match Data) และใช้วิธี Temporal Split ในการแบ่งชุดข้อมูลทดสอบ \\
\hline
\end{tabular}
\end{table}

% -----------------------------------------------------------
\section{ข้อเสนอแนะ ในการพัฒนาต่อไป}
% -----------------------------------------------------------

\subsection{การปรับปรุงฟังก์ชันการทำงาน}

\begin{itemize}
    \item \textbf{การเชื่อมต่อราคาเรียลไทม์ (Live Odds Integration):} พัฒนาระบบเชื่อมต่อ API กับเจ้ามือรับพนันเพื่อดึงราคาต่อรองล่าสุดก่อนเตะแบบอัตโนมัติ ทำให้คำนวณค่า $+EV$ ได้แม่นยำยิ่งขึ้น
    \item \textbf{การประมวลผลข่าวสารและความเคลื่อนไหว (Sentiment Analysis):} เพิ่มมอดูลวิเคราะห์ข่าวสารการบาดเจ็บ สภาพอากาศ หรือสัมภาษณ์ผู้จัดการทีมด้วยเทคนิค Natural Language Processing (NLP) เพื่อเสริมมิติการทำนาย
\end{itemize}

\subsection{แนวทางการวิจัยต่อยอด}

\begin{itemize}
    \item \textbf{การใช้ Deep Learning และ Graph Neural Networks (GNN):} ประยุกต์ใช้แบบจำลองโครงข่ายประสาทเทียมเพื่อเรียนรู้ความสัมพันธ์ระหว่างผู้เล่นในทีมแบบเชิงพื้นที่และปฏิสัมพันธ์บนสนาม
    \item \textbf{การขยายขอบเขตสู่ตลาดการเดิมพันอื่น ๆ:} ต่อยอดแบบจำลองจากการทำนายผลแพ้/ชนะ/เสมอ ไปสู่ตลาดจำนวนประตูนวม (Over/Under) หรือตลาดการต่อรองราคาแบบเอเชียนแฮนดิแคป (Asian Handicap)
\end{itemize}